% Mục 7: Các yếu tố đe dọa độ tin cậy, ba đoạn Construct/Internal/External. Mỗi giới hạn chỉ nêu ở
% đây. Nguồn: results/{apple-m5,C4M7Y670NT,LAP1184}/{campaign.log, machine.txt, lock.csv, counter.csv,
% barrier.csv, calibration.json, tsan.log}, tables.tex (campaign, lockunit, locksingle), results/apple-m5/extra/.
% Số tự tính: scripts/report_numbers/corr_stats.py, corr_stats2.py, extra_analyze.py, rlock_analyze.py.

\emph{Construct validity.} Cận dưới của phần lock dùng $t$ đo ở 1 thread, nên chứa chi phí lấy và nhả
lock của lock nhanh nhất. Ở $N=11$, chi phí này chênh 41\,\% giữa các lock
(Bảng~\ref{tab:apple-m5:locksingle}, phụ lục). Cận dưới bỏ qua giới hạn 10 core khi $p > 10$ và
$\mathrm{OW} = 15$ (ở 20 thread sàn thật cao hơn 1,6 lần). Cận dưới còn giả định mọi core chạy
\texttt{is\_prime} nhanh như ở lần chạy 1 thread. Khi số thread vượt số P-core, một phần việc chạy trên
E-core chậm hơn, nên một phần độ chậm không do lock. Ở 6 thread và $\mathrm{OW} = 15$, khi lock chưa
bận liên tục, mọi lock đã chậm 1,53--2,03 lần cận dưới.
% nguồn: locksingle N=11: 6,50/4,60 = 1,41; sàn khi p > 10: p(1+OW)It/10 so với max(p,1+OW)It, chỉ lớn hơn
% khi OW = 15 (p=12: 19,2/16 = 1,2; p=16, 20: 25,6/16 = 32/20 = 1,6); OW <= 7 thì p(1+OW)/10 <= p.
% scripts/report_numbers/rlock_analyze.py, độ chậm N=1009 OW=15 p=6: 1,53 (mcs) .. 2,03 (os_unfair_lock).
Phép quy về cùng số vòng lặp giả định thời gian tỉ lệ với số vòng. Giả định này yếu nhất ở MCS quá tải
(18--286 vòng mỗi thread trên \PCname{apple-m5}, 15--195 trên \PCname{C4M7Y670NT}, 10--56 trên
\PCname{LAP1184}), ở counter \PCname{lock_mcs} quá tải (10--200 lần tăng mỗi thread trên cả ba máy) và
ở các ô barrier sụp (50--69 episode mỗi lần chạy). Trên \PCname{LAP1184}, ở $N=1009$ và
$\mathrm{OW}=0$, phần lớn lần chạy MCS ở 16 và 24 thread xong trong vài mili giây, còn số ít mất 1,2--2,6~s.
Trung vị của hai ô này rơi vào nhóm nhanh, nên chỉ là cận dưới của chi phí MCS. Vài ô MCS quá tải chỉ
còn 4--6 giá trị sau khi vượt thời hạn (mục Phương pháp đo).
% nguồn: plot.aggregate(peers=True) trên lock.csv, mcs reduced: iters apple-m5 18..286, C4M7Y670NT 15..195,
% LAP1184 10..56 (N=10007: 10); counter.csv lock_mcs iters <= 200: apple-m5 từ 12 thread, C4M7Y670NT từ
% 16 thread, LAP1184 từ 10 thread, iters 10..200; barrier.csv (sr, sr_bo_spin, 12-20 thread: iters 69/66/50);
% LAP1184 lock.csv mcs N=1009 OW=0: p=16 5 lần 2,7..6,2 ms, 2 lần 1,18/1,73 s; p=24 6 lần 3,6..5,5 ms,
% 1 lần 2,58 s.
Ở $\mathrm{OW}=8$, thời gian của mọi counter, kể cả \PCname{faa}, tăng ít nhất gấp đôi từ 4 lên 10 thread,
và nhóm không có phép đo đối chứng để tách phần do E-core. Phép đo khoảng trống trao lock thêm hai lần
đọc đồng hồ vào mỗi critical section, nên báo cáo chỉ dùng thống kê khoảng trống của nó, không dùng thời
gian chạy. Độ công bằng so thời điểm kết thúc của các thread, nên không thấy các đợt độc chiếm luân phiên
ngắn.
% nguồn: counter.csv OW=8, giây/iters trung vị 10 thread / 4 thread: faa 2,13 (nhỏ nhất) .. lock_mcs 6,44;
% experiments/handoff_gap.cpp; định nghĩa độ công bằng: bench_lock.cpp, bench_counter.cpp (thr_min_s, thr_max_s).

\emph{Internal validity.} Không lần chạy nào ghim thread vào core. \PCname{apple-m5} có 4 P-core và 6
E-core, và macOS không cho ghim thread. Ở $N=1009$, thời gian một lời gọi \texttt{is\_prime} trong 420
lần chạy 1 thread đi từ 468 tới 1\,713~ns. Nhóm cho rằng độ phân tán này phần lớn do thread rơi vào
E-core, nhưng không tách được ảnh hưởng của E-core với tải nền. Load average 1 phút trong đợt đo lock
có trung vị khoảng 25 ở mọi số thread, kể cả lần chạy 1 thread (12,4--51,2). Con số này cao gấp đôi số
thread trung bình theo thời gian của đợt đo lock (12,0).
% nguồn: lock.csv cột load1: 12,42..51,18, trung vị 25,4 (p=1: 25,2) (corr_stats.py); số thread trung bình
% có trọng số seconds trên mọi dòng lock.csv apple-m5: 11,97 (chỉ dòng ok: 11,78).
Nhóm đo lại năm lock (\PCname{tas}, \PCname{ttas}, \PCname{mcs}, \PCname{std_mutex},
\PCname{os_unfair_lock}) ở $N = 1009$, $\mathrm{OW} = 0$, 1--10 thread, 7 lần lặp, mỗi lần chạy chỉ bắt
đầu khi load average dưới 3 hoặc sau 2 phút chờ (trung vị 2,9). Ở cấu hình này, các đường cong giữ hình
dạng khi tải thấp, còn điểm gãy của MCS ở 10 thread nhỏ đi một nửa (0,53 lần, tức chậm khoảng 7 lần
cận dưới thay vì 13). Ngoài MCS ở 10 thread, trung vị mới bằng 0,83--1,05 lần trung vị của đợt đo
chính, tức lệch tới 17\,\%. Báo cáo không diễn giải những khác biệt nhỏ hơn mức này. Lần đo lại khác
đợt đo chính cả về load average lẫn phiên chạy và khoảng nghỉ giữa các lần chạy, nên không tách được
riêng ảnh hưởng của tải. Vùng vượt số core chưa được đo lại.
% nguồn: scripts/report_numbers/extra_analyze.py (results/apple-m5/extra/lowload.csv, 175 lần chạy, load1
% 2,77..6,71, trung vị 2,89); ns/CS low/main ngoài mcs p=10: 0,83 (tas 10t) .. 1,05 (std_mutex 4t);
% mcs p=10: 3\,545 / 6\,635 ns = 0,53; t(N=1009) = 6\,635/13,0 = 510 ns -> 3\,545/510 = 6,9.
Nhóm không đọc nhiệt độ hay tần số CPU. Thời gian 1 thread không trôi quá 3,4\,\% qua 7 lượt lặp. Với 7
lần lặp, nhóm không tính khoảng tin cậy. Trung vị không thể hiện phân bố hai đỉnh, và chỉ thanh min--max
cho thấy phần này.
% nguồn: scripts/report_numbers/corr_stats2.py (lock.csv, threads=1, n=1009, giây/(iters*(1+OW)));
% scripts/report_numbers/corr_stats.py (tỉ số theo lượt 0,977..1,034)
Trên \PCname{LAP1184}, load average (5,90--19,6) chỉ đếm tiến trình trong máy ảo. Tải của Windows, việc
Windows xếp lịch cho các CPU ảo, nguồn điện và chế độ điện trong lúc đo không được ghi lại. Linux không
biết CPU ảo nào nằm trên P-core hay E-core, và \texttt{lscpu} báo tần số cố định 2\,496~MHz, nên không
thấy được throttling. Thời gian một lời gọi ở 1 thread, $N=1009$, đi từ 1\,463 tới 3\,176~ns qua 455
lần chạy. \texttt{sleep\_for} với yêu cầu 100~ns ngủ trung vị khoảng 62~\textmu s. Trong con số này,
nhóm không tách được phần của bộ hẹn giờ ảo của Hyper-V với phần của timer slack, tức khoảng trễ Linux
được phép thêm vào khi đánh thức để gộp các lần hẹn giờ. Pause-loop exiting, cơ chế hypervisor tạm dừng
một CPU ảo quay vòng chờ quá lâu, và việc đánh thức thread qua ngắt ảo là những giải thích khác, chưa
đo, cho việc barrier spin chậm ngay ở 12 thread.
% nguồn: figures/LAP1184/tables.tex (campaign); machine.txt (cpu MHz 2496.0, virtualization: Microsoft);
% lock.csv LAP1184 threads=1, n=1009: 1463..3176 ns/lời gọi, 455 lần chạy; sleep_latency.csv
% (sleep_for requested_ns=100: median_ns 61655).

\emph{External validity.} Số đo đến từ ba máy, gồm hai máy Apple Silicon (Apple clang 21, macOS) và một máy
x86-64 (g++ 11.5, Linux trong máy ảo WSL2), cả ba đều \texttt{-O2}. Hai máy Apple khác nhau ở số P-core
và số core, nhưng cũng ở thế hệ chip (M3 Pro, M5), dung lượng L2 mỗi cụm và phiên bản macOS. Vì vậy việc
điểm gãy dời theo số P-core phù hợp với giải thích hai loại core, nhưng không chứng minh nó. Máy x86-64
khác hai máy Apple cùng lúc ở kiến trúc, memory model, compiler, hệ điều hành và việc có máy ảo hay
không. Vì vậy RQ6 chỉ gán một khác biệt của máy này cho một nguyên nhân khi có bằng chứng riêng, như
\texttt{seq\_cst} với TSO. Nhóm không có máy x86 chạy Linux trực tiếp, máy chạy Windows trực tiếp, hay
máy nhiều socket. \texttt{is\_prime} không chạm dữ liệu dùng chung. Critical section thật thường ghi dữ
liệu dùng chung, nên chi phí chuyển giao lock thực tế có thể lớn hơn.
% nguồn: machine.txt apple-m5 (l2_cache 16 MiB per 4 P cores / 6 MiB per 6 E cores; macOS 27.0.1) và
% C4M7Y670NT (16 MiB per 6 P cores / 4 MiB per 6 E cores; macOS 26.7.1).
